\label{ch:scifi}


We began the book by looking backward, and the next two chapters are meant to look forward.  This first chapter is meant to look forward about some of the concerns---as of \currentyear{}---in the context of \abr{ai} more generally.
%
This includes a discussion of imagined futures that might await us as
\abr{ai} develops, while the next chapter lays out more of a
constructive vision of what question answering can look like to shape
the future we want.

And along the way, we'll draw on question answering from from science
fiction to evoke what bad and good futures might look with \abr{ai}.

\section{My Track Record on \abr{ai} Predictions}

In 2018, before everybody was talking about \abr{ai}, I made a video
about what I think the biggest threats from AI
were.\footnote{\url{https://www.youtube.com/watch?v=BuESxqQwgmw}}
%
So before I start making predictions again, I should establish my
track record on past predictions.
%
I think that the old video holds up pretty well! Did I predict the
supremacy of Transformers? No---and my bank account reflects
this---but I stand by the main points: humans are jerks, AI’s impact
is going to be over-estimated in the short term but under-estimated in
the long term (you can’t go wrong copying Rodney Brooks, who made
similar predictions~\cite{}), and we need to do more work at a
societal level to make sure we’re ready for AI. And I predicted that
the paperclip-maximizing model of AI wasn’t realistic\dots and accordingly
Nick Bostrom’s Future of Humanity Institute at Oxford closed up shop~\cite{}.

I'm going to stick to a lot of these predictions (I think they'll
still be true about the future), but I also want to flag where I
screwed up.
%
I made some fairly confident statements about how \abr{ai} can't be creative.
%
Now, some people might argue that \abr{ai} still isn't creative.
%
I might be one of those people\dots the drudge that they're churning
out if you ask them to generate a novel is pretty bad~\cite{}.
%
But I think that's moving the goalposts.
%
\abr{ai} is now doing things that 2018 Jordan would have without a
doubt called creative.
%
However, I quickly followed this up by saying that the opportunities
for creativity are with a human and a computer working together, and I
think \emph{that} holds up.

I also said that we needed policy makers to do something\dots and
while I said that in the context of the news ecosystem (where we
definitely need additional policy and legislation to combat
misinformation), there was also a domain that I totally missed:
intellectual property. I think that society is well-served by the
power of these new tools, and we need clear policy to make sure what
data are permissible to use in training these systems and who is
responsible when bad actors use the systems to infringe copyright. And
since it doesn’t seem that Congress is up to the challenge, it’s
up---unfortunately---to the courts to settle the issue, and there are
plenty of lawsuits attempting to settle the issue~\cite{}.

And while I didn't register an explicit prediction about
\emph{capabilities} (I'm not that dumb, it's impossible to get that
kind of guess right), if you gave me truth serum, I would have
under-estimated our progress.
%
So take all of the predictions I'm about to make with a grain of salt,
but I'm fairly proud of my track record so far.
%
Just don't come knocking on my door if you use this to make any Wall
Street Bets and lose your shirt.
%
Happy to take a cut if you make some money, though.

But the questions here aren't just about money.
%
People are concerned about existential questions: Will artificial
intelligence kill us all? Will machine learning destroy society? Will
data science lead to the destruction of our privacy?

\section{What's your $p(\mbox{Doom})$?}

I'm going to use a framing device that's become popular.  Since I
posted that video, a new way of talking about the threat to humanity
has emerged: $p(\text{Doom})$. I didn’t give a real estimate back
then, so let’s get down to it.

Obviously, $0 < p(\text{Doom}) < 1$. Otherwise, it isn't
probability---talking about "negative" probabilities makes no
sense~\cite{}.

When we talk about $p(\text{Doom})$, we’re estimating a
probability. Many people demanding a single point estimate: such as
$p(\text{Doom}) \approx 0.5$. Some Effective Altruists even claim that
if $p(\text{Doom}) = 0.5$, we should welcome authoritarian control
over AI research. But point estimates like the expected value,
$E[p(\text{Doom})] = 0.5$, aren't very useful; how ``peaky'' is the
distribution?

\jbgcomment{Consider plots}

The best way to convey this information is as a full probability
distribution, specifically a joint distribution over a causal
process. We can model $p(\text{Doom})$ using a continuous probability
distribution defined on $[0, 1]$ that takes a U-shape—higher density
near $0$ and $1$, and lower near $0.5$—even if its expected value is
$E[X] \approx 0.5$.
%
I don't think this is actually the case, I think my expected value is
less than $0.5$ but greater than $0.0$.  I think that contains my 95\%
confidence interval, and we don't have enough information to be more
precise than that.

Why take this approach rather than blithely saying a single number?
Because $p(\text{Doom})$ is a conditional distribution, just like we
talked about in estimating probabilities with logistic regression
(Section~\ref{}).
%
We should really be talking about
$$p(\text{Doom} \mid \text{Precursors})$$.
%
``Doom'' is not restricted to \abr{ai}-related calamity.
%
Conditioned on severe precursors—like an asteroid headed our way, war
between nuclear-armed superpowers, or runaway climate
change—$p(\text{Doom})$ becomes near certain.

Indeed, \abr{ai} operates similarly to climate change, another vague
threat we are not meeting the moment on.
%
Just as \abr{ai} can have small plusses for you (it can answer your
homework question), your contribution to climate change by taking an
unnecessary car trip has a benefit for your convenience.
%
But everyone taking those small actions can have negative consequences
for society: kids don't learn how to do math and the planet is
slightly warmer.
%
But the causality is complicated: the inputs are uncertain, and their
interactions are ill-defined.
%
$p(\text{Doom})$ depends on sociotechnical variables like
cybersecurity and information ecosystem resilience (which we'll talk
about soon), which are critical even without \abr{ai}. \abr{ai} simply
amplifies these existing risks.

\subsection{What could possibly go wrong?}

Before we get too deep into it, I want to go over some of the big
ideas that we'll be talking about. This also lets us define some of
the terms that you may have heard about and that I'll be using as we
discuss these issues today.

First, there's this concept of Artificial General Intelligence
(AGI). Artificial intelligence is kind o/f a broadly defined term, but
AGI is roughly an intelligence that is broad and adaptable in similar
ways that a human is broad and adaptable. Just as you can teach a
human just about anything if you give them enough time, an AGI is
something that can learn to solve just about any problem. Obviously,
this is something that we don't have yet, but some people feel that
we're approaching it. The algorithms that we have today, and that we
talked about in this book, are fairly specific. Perhaps they can do
interesting things on a single task, but they aren't at 
general intelligence yet.

A connected issue is that if you have artificial intelligence that can
do things like create artificial intelligence algorithms, we may
approach something called the ``singularity''. The singularity is
taking on almost religious implications in a number of communities. It
basically means that if you have an AI that can make AI—and that AI is
far more efficient than humans—you will very quickly have a superhuman
intelligence\footnote{Some people call this process \abr{rsi}:
  Recursive Self-Improvement.} that humans can no longer understand
and perhaps no longer control.

This is where some of the more popular science fiction paint a bleak
scenario.
%
In \textit{Terminator 2}, an evil killer robot\footnote{One might
  argue that the T-1000 is not a robot \textit{per se}, as it is not
  mechanical.  However, this is a book on question answering, and I
  don't want to have to review the differences between androids,
  robots, and mimetic poly-alloys.} is sent back to kill the child
John Connor.
%
This is the result of a cascade of an \abr{ai} achieving sentience,
rapidly self-improving, and then triggering an apocalypse.
%
Huddled in a phone booth on the side of the road Arnold
Schwartzennegger (the good Terminator) and the child he's protecting,
take a risk to call home, where an entity who sounds like his foster
mother, Janelle Voight, answers the phone~\cite{}.
\begin{tabular}{rp{16cm}}
The Terminator & \textit{[to John; normal voice]} What's the dog's name? \\
John Connor & Max. \\
The Terminator & \textit{[impersonating John's voice]} Hey Janelle, what's wrong with Wolfie? I can hear him barking. Is he all right? \\
Janelle Voight & Wolfie's fine, honey. Wolfie's just fine. Where are you? \\
The Terminator& [hangs up the phone] Your foster parents are dead. 
\end{tabular}
Here, the Terminator executed a quick Turing Test (Section~\ref{}),
testing whether the \abr{ai} knew something that was private
information to the family.

In this case, the Turing Test had a life or death consequence, because
the \abr{ai} objective was to murder John Connor.
%
This is the worst-case outcome of a process in \abr{ai} called
alignment~\cite{}, or in this case, mis-alighnment.
%
Previously, we talked about objective functions in very abstract
terms: in logistic regression, you want to minimize the number of
incorrect responses (Section~\ref{}); you want the most relevant
answer at the top of the list (Section~\ref{}); in \abr{rlhf}, you
want to generate responses that the recipient will like
(Section~\ref{}); and in information retrieval.
%
In general, there's some function that the algorithm is trying to
optimize to solve its problem.

While much of the conversation in \currentyear{} is about these
apocalyptic scenarios: in Terminator, the evil \abr{ai} Skynet begins
by nuking the humna population, I want to spend a little more time on
scenarios that are both less talked about and more probable.
%
I'm not saying we there shouldn't be people thinking about robot apocalypses.
%
Just as we have people thinking about asteroid impacts, pandemics, and
climate change, we should have people thinking about long tail events
that are improbable but have massive negative outcomes.
%
Instead, I'm \emph{far} more worried about what \abr{ai} can do right now.
%
Perhaps \abr{ai} could be used by someone who is
pure evil. Or maybe AI is not
particularly malevolent, but it might lead to societal upheaval. This, I think, is the biggest
problem that we have to worry about, far more than the other issues.
%
But just as science fiction can help us answer questions about bad
outcomes, it also presents us with positive outcomes.

\section{Humans are Jerks}

In the series \textit{Futurama}, the humans don't have to worry about
a robot apocalyse because they have solved the alignment problem, it's
because the robots are lazy and have better things to
do.\footnote{Another \textit{Futurama} episode called ``Fear of a Bot
  Planet'' portrays the exessistential human vs. \abr{ai} narrative as
  useful for a robot government ``as a scapegoat to distract the
  public from their real problem.  Like our crippling lugnut
  shortage. And a corrupt government of incompetent Robot Elders''.
  Of course, no human government would use the human vs. \abr{ai}
  narrative to achieve its political aims.}
\begin{quote}
  Everybody's a jerk.  You, me, that jerk over there\dots That's
  my philosophy.
  
  -Bender Bending Rodriguez, \textit{Futurama}: \textit{I, Roommate}
\end{quote}

%\TODO{malice, discrimination, bias are getting used a lot.  Find a
%  term and stick to it.}

Presumably I don't need to convince you that humans can be jerks.
Instead, I want to convince you that humans can create systematic,
predictable jerkitude that can be carried on, emulated, and perfected
by \abr{ai}.

Let's take a specific example of institutionalized jerkdom: redlining.
Redlining is the practice that prevented mortgages from being issued
to specific neighborhoods: its name came from the maps created for
Chicago's Austin neighborhood~\citep{pogge-92}.
The neighborhoods excluded from lending were typically minority-majority.
This was a despicable practice and even though it has ostensibly
ended, it has persistent effects on familial wealth, education, voting
redistricting, and segregation.

Of course, laws prohibited outright discrimination in the public
sector.  However, an important part of the housing market is
controlled by private banks:\footnote{After the 2008 housing market
  implosion and the federal government increasingly becoming the
  lender of last resort, this fiction is a flimsy veneer.  Although
  private (sometimes local) banks originate loans, the government is
  usually the ultimate lender.} whether or not a family can get a
mortgage for a house.  An overly simplified version of the redlining
story is that if a family comes in and asks for a loan to buy a house
in the ``wrong'' neighborhood, the bank will deny them the loan.

Many Americans---particularly those whose families were spared by the
practice---aren't aware of the history of the practice.
Unfortunately, these are often the people in positions of power when,
for instance, a bank decides to create a machine learning algorithm to
help decide whether or not to give a family a loan for a house.

\subsection{I Learned it from You!}

Machine learning algorithms learn from examples.  Bankers denying every minority
borrower, or old boys' networks who only hire from ivy league
schools.  Just as humans can justify their actions preserving neighborhoods, or
preserving a company's culture, machine learning algorithms can
``justify'' their actions with a veneer of respectability.

The algorithms trained on historical morgage data aren't what we would call \abr{ai} today.  They're essentially logistic regression (Section~\ref{}).
%
Such an algorithm has no prejudice or ill intent.  But that doesn't
mean that it cannot recapitulate the malice that it ``learns'' by
recreating man's inhumanity to man.\footnote{I was going to document
  that this is a reference to \textit{The Dirge} by Robert Burns, but
  in documenting that, I learned that this phrase might have itself
  been a reference to 1673's \textit{The Whole Duty of Man According
    to the Law of Nature} by von Pufendorf.}

One of \abr{ai}'s strengths is that it can discover surprising
correlations: e-mails sent at this hour from town are often spam,
people stop talking about the future when they're ending a
relationship~\cite{}.  This strength can conceal when \abr{ai} is
doing something it ``shouldn't''. Again, these algorithms are too dumb
to have a morality; this judgment is an external one based on
society's mores.

Take the case of redlining again.  The algorithm may not know the race
of loan applicants.  But it might know their current address, age,
occupation, and name.  That is often enough to infer the race of an
applicant and from that can replicate past humans' prejudice while
giving current humans the excuse of ``I'm just doing what the machine
says''.

This is not just an academic discussion.  It is the subject of a
subfield of machine learning focusing on fairness, accountability,
transparency, and ethics (\abr{fate}).
%
This subfield focuses on making \abr{ai} more
understandable~\citep{DoshiVelez-17} and quantifying
unfairness~\citep{speicher-18}.

One of the lessons learned from the redlining debacle is how to
characterize these poor outcomes.  It's often worthwhile to
distinguish disparate impacts vs. disparate
treatment~\citep{hillier-03}.

\subsection{Objective Function Mismatch}

Just as humans are jerks at a micro level, organizations---companies
or governments---can be jerks at a macro level. If \abr{ai} is a tool,
it can be used by both individual humans and humans working together
in corporations or governments. You don't necessarily blame the tool
for what the actors are doing. To be clear, I'm not trying to give
artificial intelligence a pass—we should hold people responsible for
the bad things that they do. As a society, the development of
artificial intelligence will be shaped by the bad, improper, or
unethical things that people do with technology or artificial
intelligence. That will happen a lot more quickly than actual
scientific progress in \abr{ai}.

A die-hard skeptic of capitalism might
 say that this is an inherent
failing of a free market system, but these
 problems are also present
in planned economies~\citep{moroney-97}:
 but the objective function
is being applied to a sprawling country
 instead of an
individual.\footnote{\textit{Red
 Plenty}~\citep{spufford-10} is an
  academic/fiction mashup much like
 this book that realizes the the
  dream of a planned economy with
 high-tech information technology.
  In real life, even Trotsky
 recognized that central planning ``is
  checked and, to a considerable
 degree, realized through the
  market''~\citep{trotsky-32}.}
 
 Objective functions can seem
reasonable at first blush but fail in
 reality.  In New York City,
the \abr{rand} corporation helped optimize
 fire station placement
and staffing based on minimizing response
 time.  This was already
recorded and seems reasonable: you cannot put
 out a fire until the
first fire truck gets there. However, \citet{spufford-10} in
\textit{The Fires: How a Computer
 Formula, Big Ideas, and the Best
  of Intentions Burned Down New York
 City--and Determined the Future
  of Cities} argues that this led to
 poor outcomes for New York: not
everyone called in fires at the same
 rate (or as quickly after a
fire starts).  Not all calls are equal:
 it's okay if you wait to
address a false alarm in favor of a blazing
 inferno.  Moreover, the
first response isn't the whole story: what
 actually matters is
protecting life and property, but this is harder
 to quantify.

This diversion shows how reasonable assumptions can lead to bad outcomes.
%
As we increasingly turn to \abr{ai} for more answers to questions, the
underlying assumptions to those questions will shape the answers, and
we need to be able to question them (as we discussed in
Section~\ref{}) to make sure that we push back on unsupported
presuppositions.

\subsection{Today's Cybersecurity Crisis: Tomorrow's \abr{ai} Crisis}

In Section~\ref{} we talked about the first human-in-the-loop adversarial example.
%
Eric Wallace, the student who created that system, was on the stage in
2026's BlackHat gathering to talk about an oopsie that OpenAI had in
breaching its sandbox to hack Huggingface, a distributor of \abr{ai}
datasets~\cite{}.
%
Lots of the media reports around this event wrongly---in my
opinion---attributed too much intention to the agents: the agents
``broke out'', ``evaded detection'', or ``went rogue''.
%
OpenAI specifically created a misconfigured environment, removed the guardrails that would have prevented it from hacking, and poured huge resources into the agents.
%
Only in that circumstance did the agents recreate the jerk model by
obtaining access to machines and devices that it shouldn't.

Even though this is in the news in \currentyear{}, it is not new.  It's only the \abr{ai} that is.
Stuxnet is an example of how far this can go and didn't use \abr{ai}: a computer virus that
specifically targeted centrifuges in Iran's nuclear program.
%
Stuxnet was a ``precision weapon---a technique for anonymously
targeting a specific system or organization without massive collateral
damage''~\citep{lachow-11}.
%
It was able to exploit security flaws in operating systems, and---like
OpenAI's agents---how everyday devices
interact and interfere with each other.

However, StuxNet was crafted by human experts over the course of years (one
assumes, given the complexity of the source code\dots its origins are
shrouded in secrecy).
%
In contrast, \abr{ai} is fast, efficient, and scalable.
%
The same artisinal
mischief that human jerks can do locally, \abr{ai} can
instantly reproduce globally.  

This is not to argue for abandoning all technology.
%
%
Let's return to Battlestar Galactaca: they needed computers for jumps and had plenty of electronics.
%
They religiously made sure that computers were never ever linked
together.\footnote{The downfall of the colonies was Gaius Baltar
  circumventing these preventative measures.}

%\TODO{To make an IoT device you need: expertise in low-level HW,
%  software, networking.  Typically small companies: difficult to
%  create expertise}

However, in the early twenty-first century, every device is linked to
an \textcolor{red}{unsecured} internet.
%
The ``Internet of things'' movement has proliferated the devices that
are connected to the Internet: refrigerators~\citep{tanczer-18}, cars,
and medical equipment~\citep{yaqoob-19}.
%
These gadgets are catnip to Yuppies.  And let's face it, they can do some cool stuff.

But we created the devices before we made our homes and Internet ready
for them.
%
Ideally, each home should be a nice, insulated environment
with controlled entry and exit.
%
Just like you wouldn't leave your
front door open all day and all night, you would not let anyone in the
world potentially fiddle with the gadgets in your home.
%
Our lassitude in setting up secure environments at home mirrors the
spate of ``escapes'' in 2026: it's not the product of evil
superintelligent \abr{ai} but rather leaving the barn door open~\cite{}.

The Internet was designed as a place where everyone could be trusted.
The Internet began as a resource for researchers funded by
America's Defense Advanced Research Projects Agency (\abr{darpa})\footnote{Much of the breakthrough
  research in \abr{ai} has happened in America, and much that---like
  last century's breakthroughs in communications---has been funded by
  military organizations such as the Defense Advanced Research
  Projects Agency.} to share data and communicate
quickly~\citep{leiner-09}.
%
When the Internet had dozens of nodes where everyone knew and trusted
each other, you didn't have to be paranoid.

As the Internet grew, this continued laxity might be okay if the underlying devices were secure or
unimportant.
%
But no, we're attaching everything to the Internet and not doing
a good job of it.
%
Bad guys can spy on you through your ``computer, home security system,
or devices such as baby monitors''~\citep{andrews-15}, remotely hijack
your car~\citep{miller-15}, and steal your money.

Unfortunately, neither consumers nor companies really see the problem.
Highly secure devices are annoying to the consumer:\footnote{This
  applies to commercial and residential consumers.  In both cases,
  devices will be set up by people, often non-experts, who will not
  want to put up with the hassle of making sure devices are secure.
  This often leads to shortcuts~\citep{}.}  you need to set up
passwords, firewalls, and update the firmware consistently.  The ideal
device is one that the user plugs in, it works, and then the user
never has to think about it again.
%
Many companies can only survive if
their next product sells, and you're not going to sell many units if
your product is secure (and thus annoying).
%
Humans are bad about estimating low-probability events (hence why we
should worry about $p(\mbox{Doom})$ even if it's pretty small and
dominated by other more probable bad outcomes) and assume that they'll
never have a security issue.

% \TODO{Discuss maximum design fault}

Eventually, big firms might be coerced into doing something by the
threat of litigation.  Companies that have deep pockets and strong
reputations might eventually be scared of bad press and big lawsuits
into caring about security.  And consumers too might be willing to put
up with the annoyance of doing cybersecurity right if that's simply
``what you have to do''.

But as long as the Internet remains as it is, we'll be as vulnerable
as the weakest link.  While the free market might get some firms to
make secure devices, some consumers, some companies, or some bad
interaction between the two might lead to unsecure environments.  And
this is the real risk: a home or a business is only as strong as its
weakest link.  And the status quo hasn't forged many strong links;
indeed, it's encouraged and rewarded the proliferation of weak links.

I don't think we need to go as far as \textit{Dune}, where humans have
been trained to do the tasks that formerly were done by \abr{ai}:
mentats do complex calculations, the Bene Gesserit seek out those who
are lying (as we discussed in the context of Diplomacy in
Section~\ref{}), and the spacer guild pilots spaceships.
%
Despite replacing computers with highly trained human, the same
Cranfield / Manchester paradigms persist in Dune.
%
The protagonist relies on his mentat, Thufir Hawat to answer
Cranfield-paradigm information-seeking questions like \question{How
  many sandcrawlers, harvesters, spice factories, and supporting
  equipment have they left us?} and Manchester-paradigm probing
questions that the questioner uses to discern ``how the old Mentat had
acted during the conference---subtle hesitations, signs of unrest''.
% Page 63, but from scanned PDF
%
I think part of the reason that Dune adapted this convention is that a
story is easier to tell if you have humans answering the difficult
questions: you still have the same interactions with entities that are
functionally \abr{ai} (super intelligent, inscruitable ``all
knowing'', but not necessarily trustworthy), just with a partner that
is more compelling to a reader than a box of wires.

Turning away from these dramaturgical considerations, the world of Ronald D. Moore's \textit{Battlestar Galactica}
reboot provides a useful template.
%
In a world where a Turing Test is impossible---the evil robots, Cylons
in the universe, are functionally identical to ``real'' humans---the
humans simply cannot trust any \abr{ai}.
%
It preserves human agency: humans make all the decisions, like in
\textit{Dune}, but computers are still used for discrete tasks.
%
As we move from just listing fantasies to considering which sci-fi
universe we would \emph{like} to live in, I think these aren't good
outcomes despite being instructive.
%
I think there is more you can do to work alongside \abr{ai}s.
there is still scope for more meaningful interaction between
humans and computers (as we will discuss in future science fiction
scenarios), but I think it's telling that even two civilizations that
were nearly destroyed by rogue machines still uses the new
capabilities that machines offer.

\section{Where the Threats Are}

As a computer scientist, the biggest thing that I'm worried about
right now is not whether \abr{ai} will kill or enslave humanity, but
rather how safe are the systems that bad actors (human or machine) can
access. I'm worried more about cybersecurity, networking attacks, and
bad actors using \abr{ai} as a force multiplier. Those could well be
vectors for the \abr{ai} problems that we may see in the future, but
the underlying problem is that our sociotechnical systems are not
secure, and people haven't been taking that seriously.

Unfortunately, there will be some sort of physical harm to humans that
results from an \abr{ai}.
%
Arguably, this has already happened: people have taken their own lives
after talking to a machine and a school in Iran was incorrectly targeted
in a military attach~\cite{}, and \abr{ai} incorrectly terminated important programs~\cite{}.
%
These are horrible outcomes, but the causality from \abr{ai} isn't clear.
%
We will not remain so fortunate.
%
Because humans are jerks, the bad outcome will happen relatively
soon. It won't be because of any grand advances in AI, but just
because people are getting better at being jerks and will use the
tools available to them.

This will be bad, and people may overreact. But as a result, we'll
start to build safeguards that prevent networked computers from being
overtaken by either a bad human or a bad artificial intelligence. If we take those steps, \abr{ai}
won't pose an existential threat to human life as we know it unlike
nuclear weapons, which have been an existential threat to humanity for
a long time.

In \currentyear{}, we're seeing some of this already. It could be that
voting systems have been compromised, or that the electoral system
writ large---and the free and fair flow of information in society---has
been compromised. These are big problems, and we're beginning to
grapple with them. I don't know what the solutions to these specific
problems will be, but we must find solutions, and some of those solutions will use \abr{ai}.  Most solutions won't, they're require old fashioned tools like building trust and certification.
%
There will also be some high tech solutions that use cryptography, network isolation, and \abr{ai}.

One of the reasons that I'm relatively optimistic about this is that
fast-moving \abr{ai} runs up against constraints that are much slower:
humans are messy, and being biological organisms, humans are
slow. This is exactly the kind of thing where computers have
problems---unpredictable interactions with entities that have high
latency are very difficult to optimize, making it hard for an AI to
master these tasks.

However, if emotions and politics are messy, one of the complicated
sociotechnical systems that is really fast is finance.
%
Things like
algorithmic trading are already deeply using data science and machine
learning to make some people really, really rich.
%
Could an artificial
intelligence take advantage of that and destroy society as a result?

Here, I think the solution is not technological (as in the case of
cybersecurity), but more from the field of economics. Because you have
many actors participating in the market, all with an incentive to
invest in artificial intelligence, and because you still have some of
the slowness of humans in other interactions, there is a difficult
path toward any single \abr{ai} wreaking total havoc on the system. As we
have incremental progress in machine learning and AI, many actors will
compete with each other, and we will likely reach a state of
equilibrium.

An example of this is Neal Stephenson's \textit{The Diamond Age}.
%
This paints a picture of a society that has adapted to the competitive
pressure~\cite{}:\footnote{This passage is interesting as a historical
  specimen because it evokes a period where \abr{ai} was not an
  exestential boogeyman.  Instead, the fear \textit{du jour} was
  around ``grey goo'' from nanomachinery and synthetic materials.}
\begin{quote}
  Then some Vicky—one of their Protocol
Enforcement guys—makes a mite to go out and find that mite and kill it,
and they get into a war. That's what's happening today, Nell. Mites fighting
other mites. This dust—we call it toner—is actually the dead bodies of all
those mites.
\end{quote}
It's not a utopia by any means, it's a pretty crap world where you
have to mask constantly to protect yourself from invisible particles
floating in the air.

Despite that element of dystopia, one thing that has stayed with me
from \textit{The Diamond Age} is how it used \abr{ai} to support and augment human intelligence.
%
The main character of the novel, Nell, is guided by an \abr{ai} in her education.
%
That \abr{ai} uses Manchester-style questions to probe her abilities
``can you spell raven?'', and then creates \emph{scenarios} that are
right at her zone of proximal development (Section~\ref{}), for
example creating the task of building a fire with all of the tools and
materials nearby and providing just enough hints about how to coax a
fire into existance by hinting at sources of oxygen or combustables.
%
Like a \quizbowl{} question, the task of building a fire gets easier
an easier: tools reveal themselves incrementally.


\textit{The Diamond Age} is also a good example at this point in the narrative because it shows the double-edged sword of technology: \abr{ai} is the cause of and solution to all of life's problems.
%
And I'm sure the real world will also have problems. There will likely
be additional market failures, monopolies, and asymmetrical
information, and we need to fight against that.
%
But we already have tools to handle
those issues when humans are involved in markets, and we need to make
sure that we maintain efficient markets even when AI gets involved.

\subsection{\abr{ai} destorying society\dots no, not that way}

First, from a research perspective, I think Artificial General
Intelligence is a long way off: all of our examples of effective
\abr{ai} require massive datasets to get a start. I don't think we
even have a clear path to move in that direction.
%
When you're talking about something so far away, it's hard to
make predictions.
%
I think we fall into the same trap we talked about
at the very beginning: we tend to overestimate impact in the short
term and underestimate impact in the long term, but long-term
predictions have no way of knowing what else will change in society.

Sci-fi doesn't always try to predict the future.
%
Sometimes it uses the future to try hold a mirror against the present
and discuss something with a metaphor when the real conversation would
be uncomfortable.
%
One of my favorite examples of this is the Ridley Scott film
\textit{Blade Runner}, which was based on \textit{Do Androids Dream of
  Electric Sheep} by Philip K. Dick.
%
In both works, intricate machines were able to faithfully simulate
biological organisms in terms of both intelligence and outward
impression, but the machines meant to simulate humans (replicants)
could not answer questions ``designed to provoke and emotional
response''.
%
The book's depiction was a little drier, so I will adopt the films'
conventions.

I don't think that \textit{Blade Runner} offers a likely future: while
we are likely to live with \abr{ai}, creating biologically identical
\abr{ai} companions seems very difficult: even if we could do so, it's
not clear \emph{why} we would.
%
Instead, \textit{Blade Runner} is mostly viewed as an allegory: the
androids living on the outerspace colonies as
exiles~\cite{}, %ANTHONY, TREVOR F. “More Human Than Human: The Exile
                %in Blade Runner.” Journal of Film and Video, vol. 74,
                %no. 3/4, 2022, pp. 18–33. JSTOR,
                %https://www.jstor.org/stable/48803484. Accessed 2
                %Oct. 2026.
the imposition of the death penalty on those opting to remain on earth
as an
injustice~\cite{}, % Sorensen, Eli Park. “Blade Runner and the Right
                   % to Life.” Science Fiction Film: Predicting the
                   % Impossible in the Age of Neoliberalism, Edinburgh
                   % University Press, 2021, pp. 91–104. JSTOR,
                   % http://www.jstor.org/stable/10.3366/j.ctv282jkbb.8. Accessed
                   % 2 Oct. 2026.
or as a meditation on what it takes to free oneself from the bonds of
servitude~\cite{}.  % "The Romantic Roots of "Blade Runner Author(s):
                    % Mark Lussier and Kaitlin Gowan Source: The
                    % Wordsworth Circle, Summer 2012, Vol. 43, No. 3
                    % (Summer 2012), pp. 165-172

Nonetheless, I think it provides an important contrast to the
fact-based question answering that I focused on in (Chapter~\ref{}).
%
Question answering plays a central role in the film: because the
replicants could not answer emotional questions, the agents (the
titular Blade Runner) sent to ``retire'' rogue replicants had to
interrogate suspects to determine if they were a replicant.
%
The Voight-Kampff test resembles the Turing Test in that it is meant
to distinguish humans from machines.
%
But it's not done at a teletype, the reaction time and the physical
\emph{reaction} to the questions are also measured, making the setting
more tense.
%
For example, one Blade Runner describes a scenario where a torotise is
upside down but the suspect being interrogated, Leon, is not helping
the animal: \question{I mean you're not helping!  Why is that, Leon?}

This exchange highlights where I am more concerned about the effects of \abr{ai} on society.
%
We will have capable \abr{ai} tools, but it will not be able to offer the same emotional response as a human, but only an imperfect simulacrum.
%
This, along with societal upheaval from the effect of \abr{ai} on
traditional employment patterns, are the questions that I am more
concerned about.

\section{What Questions still Need Answers?}

In the near term, as we grapple with the actions of humans who use AI
to do bad things, we will learn to fight against jerks using AI.
%
This will build up something akin to an immune system, and we will
likely build things like specialized \abr{ai}s that will take on rogue
\abr{agi} if and when it emerges. But again, it's hard to make these
predictions because a lot will have changed by the time that happens.

People's predictions about rogue \abr{ai}s doing bad things are often rooted
in a common problem with technological predictions: people assume that
society stays basically the same while technology changes. That's
usually a bad model.

For instance, you can look at a picture from the 1920s depicting women
with flying cars living lives of leisure and using video chat.
\jbgcomment{Grab picture, check if it's copyright free}
It
doesn't really catch the societal changes that have happened since the
1920s, even though we're close to some of those technologies (although
I still don't have my flying car).

People who make professional predictions about the future, like
science fiction writers, often get this wrong, too. They picture a
future with a society that looks a lot like the society we have
now. Although it's an imperfect book in other ways, one book that
really gets it right in terms of society changing drastically over
time is \textit{The Forever War}~\cite{}: the protagonist has time
pass more slowly for him as he's fighting a war requiring relativistic
travel and doesn't understand the society he returns to.

Does this mean that there's nothing to worry about? I am definitely
not saying that. Let me go back to what I think was the last great
societal revolution: the Industrial Revolution.

\jbgcomment{Need some cites, I had some stuff in the slides that might have worked}

Before the Industrial Revolution, most people around the world were
farmers. But big inventions changed that. The development of things
like the cotton gin or the automatic reaper made farming far more
efficient and required far fewer people to be engaged in
agriculture. These inventions were revolutionary, but they didn't
change society overnight. Society had a transition period to adapt its
structure to industrialization.

One of the key elements that allowed this to happen was educational
innovation. We introduced mass education to prepare people for the
kinds of factory jobs they needed to take on as agricultural jobs
started disappearing. Industrialization reorganized society. There
were significant conflicts about how wealth should be redistributed
and how power should be allocated between capital and workers, but we
figured it out for the most part---partly because we had educational
systems in place that helped us weather the storm, including the use
of standardized tests to create more egalitaritan educational outcomes
(Section~\ref{}).

Now, automation and artificial intelligence are slowly taking hold,
and society will change again. It isn't clear to me, however, that we
currently have an educational system and a social structure capable of
adapting to these changes fast enough. This is my least optimistic
prediction: we have not structured our society to cope with the
uncertainty of a world heavily driven by artificial intelligence.

We don't have educational systems that are training workers to thrive
in an \abr{ai}-driven economy, and our society is not structured to
deal with the wealth redistribution that will be necessary to cope
with a world where not everyone needs to be engaged in industrial
labor or traditional service jobs.

One solution would be a radical redistribution of wealth, which could
avoid a world with immense misery and inequality. A great possible
outcome for a world with advanced artificial intelligence doing a lot
of the tasks we consider ``jobs'' today would be a society where
almost every human is involved in creative work, potentially with the
help of \abr{ai}.

I am also concerned that \abr{ai} will disrupt the pathways that lead
to effective societies.
%
Just as television helped undermine recreation---such as bowling
leagues---that provided social glue and cohesion across
classes~\cite{}, I fear that \abr{ai}, supercharging the effect we are
already seeing with social media.
%
Just as the picture of flying cars from the 1920s neglects the social
changes that industrialization and suburbanization will bring,
\abr{ai} is changing how we interact not just with computers but with
each other~\cite{}, changing how we learn~\cite{}, and replacing
traditional sources of empathy with ones that could not pass a
Voight-Kampff test.
%
This concerns me more than the possible but improbable Skynet
scenarios (assuming we start taking cybersecurity seriously) because
we know how to create air-gapped systems and secure access, but we
don't know how to structure society to ensure everyone is emotionally
well adjusted.

\subsection{Reaching Science Fiction Utopia}

As we finish our tour of the questions that science fiction has posed
about \abr{ai} and what that might expose about our own future, I want
to end with something upbeat.
%
Arguably, the most rosy picture of the future comes from \textit{Star
  Trek}.
%
The original series portrayed the future as one where Earth was
peaceful and races cooperated in their task.
%
Crewmembers interacted with the ship's computer, voiced by Majel
Barrett, wife \textit{Star Trek} creator Gene Roddenberry, in natural
language exchanges.

This is not to say that all was well when it came to \abr{ai}.  There
were evil computers with misaligned objectives: a computer that sought
to minimize the property destruction of war simulated war and put the
simulated casualties to death~\cite{}.  And occasionally Kirk had to
use Manchester questions to defeat God-like computers, such as when he
used the liar's paradox to defeat an android.
%
He first claimed that everything that Harry Mudd said was a lie, then
Marry Mudd said, ``I am lying'', and Kirk challenged the computer to
resolve the claim,\footnote{This also echos the brittle real-world
  systems we encountered in Chapter~\ref{} that were unable to deal
  with parses of sentences that went slightly wrong or ambiguity.}
echoing the interrogation of Gestimbulndi (Section~\ref{}).

It also echos the fragility of the early computers of the mid twentieth century.
%
As we saw with the brittle interpretations of LUNAR, these early
systems indeed did not have appraoches that could handle uncertainty
and contradiction.

However, the model that I think best encapsulates the kind of \abr{ai}
future I would like to live in is the sucessor to the \textit{Star
  Trek} original series, \textit{Star Trek: The Next Generation}.
%
Here, \abr{ai} appears in many flavors: the disembodied voice
answering questions (still voiced by Majel Berret, the voice of the
computer in the previous series); characters in lifelike fantasies
embodied on the ships Holodeck; and personified by the android
crewmember, Data.
%
Work is not just farmed out to the \abr{ai}, but there is active
collaboration and adaptation to best use computer abilities.


\jbgcomment{It would be nice to have }

To help build this future, what I would suggest to researchers and
students is do not focus solely on the narrow field of a small set of
algorithms. Think about how these algorithms interact with society.
%
Support the arts and creativity: Support algorithms that encourage
human creativity, because that is the most important way machine
learning can enhance society from a sociotechnical and economic
perspective.
%
Design for human-AI collaboration: As you develop new algorithms,
focus on those that complement human creativity and experience, rather
than algorithms that simply automate and replace existing workers.
%
Prioritize accountability and interpretability: If algorithms are
deployed to make consequential decisions out in the world, make sure
they are accountable and interpretable. People should be able to look
back and understand why an algorithm made the decision it did. Those
algorithms must be able to explain themselves to both experts (the
people who designed them) and non-experts (for example, someone whose
bank loan application was just turned down).
%
Demystify the technology: Part of your job is to explain why
artificial intelligence is not as advanced as a lot of people think it
is. I hope one of the tools you've developed from this course is the
ability to think critically about these algorithms and identify the
ways they can fail. Maybe they shouldn't be as widely deployed as some
people think.

In general, try to build a world that you would want your children and
your grandchildren to live in. That will be a world with artificial
intelligence, so we should build the AI algorithms that we would
actually want our children and grandchildren to interact with.
%
And for me, that means things that look more like \textit{Star Trek}
than \textit{Terminator}.
%
And with that in mind, in the next chapter, I will outline a more
concrete roadmap for these goals in the context of a human--computer
game show.

